% -----------------------------------------------------------------------
%  lifelines-hc Application Note — Bioinformatics journal
%  BULLET-POINT OUTLINE VERSION  (for review / author editing)
%
%  Identical structure and front matter to main.tex.
%  All narrative paragraphs replaced with bullet-point lists.
%  Equations, code listing, figure, and bibliography unchanged.
% -----------------------------------------------------------------------
\documentclass[unnumsec,webpdf,contemporary,large]{oup-authoring-template}

% ---- packages ----
\usepackage{microtype}
\usepackage{booktabs}
\usepackage{listings}
\usepackage{xcolor}
% tighter itemize spacing without enumitem (which conflicts with OUP class)
\newlength{\savedtopsep}
\let\origitemize\itemize
\let\origenditemize\enditemize
\renewenvironment{itemize}{%
  \origitemize
  \setlength{\itemsep}{1pt}%
  \setlength{\parskip}{0pt}%
  \setlength{\topsep}{3pt}%
}{\origenditemize}

\lstset{%
  language=Python,
  basicstyle=\ttfamily\small,
  keywordstyle=\color{blue!65!black}\bfseries,
  commentstyle=\color{gray!70},
  stringstyle=\color{green!45!black},
  breaklines=true,
  frame=tb,
  framesep=4pt,
  backgroundcolor=\color{gray!7},
  xleftmargin=1.2em,
  xrightmargin=1.2em,
  aboveskip=0.8em,
  belowskip=0.5em,
  showstringspaces=false,
}

\graphicspath{{../figs/}}

% -----------------------------------------------------------------------
\begin{document}

\journaltitle{Bioinformatics}
\DOI{DOI HERE}
\copyrightyear{2025}
\pubyear{2025}
\access{Advance Access Publication Date: Day Month Year}
\appnotes{Application Note}

\firstpage{1}

% ---- title ----
\title[{\texttt{lifelines-hc}}: Higher Criticism survival testing]%
      {{\texttt{lifelines-hc}}: a Python package for HCHG testing
       in two-sample survival analysis}

% ---- authors ----
\author[1,$\ast$]{Alon Kipnis}
\author[2]{Zohar Yakhini}

\authormark{Kipnis and Yakhini}

\address[1]{%
  \orgdiv{School of Computer Science},
  \orgname{Reichman University},
  \orgaddress{%
    \street{Herzliya},
    \postcode{4610101},
    \country{Israel}%
  }%
}

\corresp[$\ast$]{%
  E-mail: \href{mailto:alon.kipnis@runi.ac.il}{alon.kipnis@runi.ac.il}%
}

\received{}{0}{}
\revised{}{0}{}
\accepted{}{0}{}

% ---- structured abstract ----
\abstract{%
\textbf{Motivation:}
\begin{itemize}
  \item Log-rank has good power against many alternatives, but
        \citet{kipnis_biometrika} showed it is \emph{not} optimal
        against sparse hazard departures — concentrated differences
        in a few intervals at \emph{unknown} locations.
  \item Weighted log-rank alternatives (Gehan-Wilcoxon, Tarone-Ware,
        Peto-Prentice, Fleming-Harrington) each pre-specify a temporal
        emphasis and are insensitive to sparse departures outside that
        window.
\end{itemize}
\textbf{Results:}
\begin{itemize}
  \item \texttt{lifelines-hc}: Python package extending \texttt{lifelines}
        with the HCHG test (HC statistic applied to per-interval
        hypergeometric p-values; \citealt{kipnis_biometrika}).
  \item HCHG detects sparse hazard departures at unknown locations with
        no pre-committed temporal weight.
  \item Four real clinical datasets: CheckMate~057 PFS ($n=582$),
        COMET-1 OS ($n=1028$), AZURE DFS ($n=3359$), and the CSL
        liver-cirrhosis trial ($n=446$) — the last with fully public
        individual patient data.
  \item HCHG: $p \leq 0.014$ in all four. Log-rank non-significant
        throughout ($p \geq 0.26$); among the weighted tests only
        Gehan-Wilcoxon and Peto-Prentice on CheckMate~057 reach $p<0.05$.
\end{itemize}
\textbf{Availability and Implementation:}
\begin{itemize}
  \item \url{https://github.com/alonkipnis/lifelines-hc}; MIT licence.
  \item \texttt{pip install lifelines-hc};
        requires Python~${\geq}\,3.8$, \texttt{lifelines}~${\geq}\,0.29$.
  \item Full analysis scripts and per-method results tables in
        Supplementary Data.
\end{itemize}%
}

\keywords{survival analysis; higher criticism; sparse signal detection;
  hypergeometric test; hypothesis testing; Python}

\maketitle


% =======================================================================
\section{Introduction}
% =======================================================================

\textbf{Background: log-rank and its limitation}
\begin{itemize}
  \item Kaplan-Meier \citep{kaplan1958} and log-rank \citep{mantel1966}
        are the standard tools for two-sample survival comparison.
  \item \citet{schoenfeld1981}: log-rank is asymptotically optimal
        against proportional-hazards (PH) alternatives.
  \item \citet{kipnis_biometrika}: log-rank is \emph{not} optimal
        against \textbf{sparse} hazard departures — hazard differences
        confined to a small, \emph{unknown} subset of time intervals.
\end{itemize}

\textbf{Motivating biological examples of sparse departures}
\begin{itemize}
  \item \textit{Crossing curves (immuno-oncology):} PD-1 inhibitors
        require an immune-priming phase; hazard excess is confined to
        early and late phases. CheckMate~057 (nivolumab vs docetaxel,
        NSCLC) \citep{borghaei2015}.
  \item \textit{Transient early benefit (targeted therapy):} Cabozantinib
        (MET/VEGFR2 inhibitor) improved bone scan response at week~12 in
        COMET-1 (mCRPC) but the OS advantage was confined to the
        bone-response phase \citep{smith2016}.
  \item \textit{Subgroup-dependent accumulated benefit (pharmacology):}
        Zoledronic acid's bone-microenvironment benefit is
        menopause-dependent; hazard difference accumulates only after
        premenopausal patients transition to postmenopause.
        AZURE trial \citep{coleman2011}.
  \item \textit{Time-varying treatment effect (hepatology):} Prednisone
        in liver cirrhosis is harmful early (infection / fluid retention)
        but anti-inflammatory later; the survival curves cross more than
        once, with the effect concentrated in a few non-contiguous windows.
        CSL trial \citep{csl_data} --- a canonical non-PH dataset with
        \emph{fully public} individual patient data.
\end{itemize}

\textbf{Existing non-PH alternatives and their limitation}
\begin{itemize}
  \item Weighted log-rank family: Gehan-Wilcoxon \citep{gehan1965},
        Peto-Prentice \citep{peto1972}, Tarone-Ware \citep{taroneware1977},
        Fleming-Harrington \citep{flemingharrington1991}.
  \item Each applies a fixed temporal weight (early / mid / late emphasis).
  \item Requires knowing \emph{a priori} where the hazard departure will
        occur — not suitable for the unknown-location sparse setting.
  \item All four can simultaneously miss a sparse departure if its location
        falls outside the intended window (demonstrated below).
\end{itemize}

\textbf{The HCHG test and the present work}
\begin{itemize}
  \item \citet{kipnis_biometrika} proposed HCHG: the HC statistic of
        \citet{donoho2004} applied to per-interval HyperGeometric (HG)
        p-values from survival data.
  \item Designed specifically for sparse departures at unknown locations;
        no pre-specified temporal weight.
  \item No accessible Python implementation existed.
  \item \textbf{This paper:} \texttt{lifelines-hc} extends
        \texttt{lifelines} \citep{davidsonpilon2019} with HCHG and
        supporting diagnostics.
\end{itemize}


% =======================================================================
\section{The \texttt{lifelines-hc} package}
% =======================================================================

\subsection*{The HCHG statistic}

\textbf{Step 1 — Hypergeometric interval p-values}
\begin{itemize}
  \item Partition follow-up into $K$ equal-width intervals.
  \item Interval $k$: $n_k$ total events, $r_k^A$/$r_k^B$ at risk,
        $x_k$ events in group B.
  \item Under the null of equal group-specific hazards:
\end{itemize}
%
\begin{equation}\label{eq:hypergeom}
  x_k \;\Big|\; n_k,\, r_k^A,\, r_k^B
  \;\sim\;
  \mathrm{Hypergeometric}\!\left(r_k^A + r_k^B,\; r_k^B,\; n_k\right),
\end{equation}
%
yielding a one-sided p-value $p_k$ per interval.

\textbf{Step 2 — HC aggregation (\citealt{donoho2004})}
\begin{itemize}
  \item Order: $p_{(1)} \leq p_{(2)} \leq \cdots \leq p_{(K)}$.
  \item Apply the HC statistic to the ordered p-values:
\end{itemize}
%
\begin{equation}\label{eq:hc}
  \mathrm{HCHG}_K
  \;=\;
  \max_{1 \leq j \leq \lfloor\gamma K\rfloor}
  \frac{j/K - p_{(j)}}{\sqrt{p_{(j)}\!\left(1 - p_{(j)}\right)/K}},
\end{equation}
%
\begin{itemize}
  \item $\gamma \in (0,1)$, default $0.2$: restricts scanning to the
        most extreme fraction of p-values \citep{donoho2004}.
  \item Global p-value: label permutation ($n_\mathrm{perms}=2000$
        in the analyses reported here).
  \item Two-sided: use $\min$ of two one-sided $p_k$ in place of $p_k$.
\end{itemize}


\subsection*{Software design and API}

\textbf{Main functions}
\begin{itemize}
  \item \texttt{higher\_criticism\_test()} — HCHG test; returns
        \texttt{StatisticalResult} with \texttt{p\_value} and
        \texttt{test\_statistic}.
  \item \texttt{fisher\_combination\_test()} — Fisher combination
        ($-2\sum_k \log p_k$) \citep{fisher1932}.
  \item \texttt{min\_p\_test()} — Bonferroni-corrected MinP test.
  \item \texttt{suspected\_deviations()} — per-interval HG p-values
        with \texttt{suspected} boolean column (HCHG-flagged intervals);
        used for diagnostic KM shading.
  \item \texttt{run\_all\_tests()} — benchmarks HCHG, Fisher, MinP,
        log-rank, and four weighted alternatives; returns
        \texttt{pandas.DataFrame}.
  \item Plotting: \texttt{plot\_km\_with\_hc()},
        \texttt{plot\_pvalue\_profile()}.
\end{itemize}

\textbf{Usage example}
\begin{lstlisting}
from lifelines_hc import (higher_criticism_test,
                           fisher_combination_test,
                           min_p_test,
                           suspected_deviations)

result = higher_criticism_test(
    T_A, T_B,
    event_observed_A=E_A, event_observed_B=E_B,
    n_intervals_to_pool=100, n_permutations=2000,
    alternative='both', gamma=0.2, seed=42)

print(result.p_value)       # permutation-calibrated p-value
print(result.test_statistic)  # HCHG statistic value
\end{lstlisting}

\textbf{Design notes}
\begin{itemize}
  \item API consistent with \texttt{lifelines.statistics}.
  \item All three omnibus tests share the same HG interval p-values and
        permutation infrastructure.
\end{itemize}


% =======================================================================
\section{Results}
% =======================================================================

\textbf{Overview}
\begin{itemize}
  \item Four real clinical datasets with distinct sparse non-PH patterns
        (Figure~\ref{fig:main}). Top row (A--D): KM curves;
        bottom row (E--H): per-interval $p$-value profiles.
  \item Benchmark: HCHG vs log-rank vs four weighted alternatives.
  \item Three datasets use IPD reconstructed from published KM curves;
        the CSL trial (Domain 4) has \emph{fully public} raw IPD.
  \item Full per-method p-value tables in Supplementary Data.
\end{itemize}

\subsection*{Domain 1 — Immuno-oncology: crossing survival curves}

\textbf{Trial and data}
\begin{itemize}
  \item CheckMate~057 \citep{borghaei2015}: nivolumab ($n=292$) vs
        docetaxel ($n=290$), 2nd-line advanced non-squamous NSCLC.
  \item Endpoint: progression-free survival (PFS).
  \item IPD reconstructed via \citet{guyot2012} / \texttt{kmdata}.
\end{itemize}

\textbf{Survival pattern}
\begin{itemize}
  \item Curves \emph{cross}: docetaxel has an early advantage (immune
        priming), then nivolumab gains a durable late benefit.
  \item Sparse departure: elevated hazard for nivolumab confined to the
        early crossing window; then a late hazard reversal.
\end{itemize}

\textbf{Test results}
\begin{itemize}
  \item Log-rank: $p = 0.351$ (NS) — early excess and late benefit cancel.
  \item Gehan-Wilcoxon: $p = 0.041$; Peto-Prentice: $p = 0.049$
        (borderline) — their early emphasis happens to align with the
        crossing; not a general signal.
  \item Tarone-Ware: $p = 0.358$; Fleming-Harrington\,(1,1): $p = 0.256$
        (both NS).
  \item HCHG: $p = 0.001$ — flags both the early and late divergence
        windows independently (Figure~\ref{fig:main}A,E).
\end{itemize}


\subsection*{Domain 2 — Targeted therapy: transient early benefit}

\textbf{Trial and data}
\begin{itemize}
  \item COMET-1 \citep{smith2016}: cabozantinib ($n=682$) vs prednisone
        ($n=346$), chemotherapy-pretreated mCRPC. 2:1 randomisation.
  \item Endpoint: overall survival (OS). Median OS $\approx 10$ months.
  \item IPD reconstructed via \citet{guyot2012} / \texttt{kmdata}.
\end{itemize}

\textbf{Biological mechanism and survival pattern}
\begin{itemize}
  \item Cabozantinib inhibits MET and VEGFR2: reduces bone lesion
        activity (bone scan response at week~12 significantly improved).
  \item Resistance through alternative pathways develops rapidly;
        benefit is confined to the bone-response phase (months 3--7).
  \item Sparse departure: early OS advantage, then curves converge.
\end{itemize}

\textbf{Test results}
\begin{itemize}
  \item Log-rank: $p = 0.262$ (NS).
  \item All four weighted alternatives: $p \geq 0.19$ (NS) — early-weighted
        tests (GW, PP) partially align but are diluted by the post-response
        null period.
  \item HCHG: $p = 0.014$; Fisher combination: $p = 0.028$
        (Figure~\ref{fig:main}B,F).
\end{itemize}


\subsection*{Domain 3 — Adjuvant bisphosphonate: menopause-dependent effect}

\textbf{Trial and data}
\begin{itemize}
  \item AZURE \citep{coleman2011,coleman2014}: zoledronic acid ($n=1681$)
        vs control ($n=1678$), early-stage breast cancer. 10-year follow-up.
  \item Endpoint: disease-free survival (DFS). 973 events total.
  \item IPD reconstructed via \citet{guyot2012} / \texttt{kmdata}.
\end{itemize}

\textbf{Biological mechanism and survival pattern}
\begin{itemize}
  \item Zoledronic acid targets the bone microenvironment (osteoclast
        inhibition), preventing micrometastasis in the bone niche.
  \item Benefit is menopause-dependent: requires low oestrogen /
        high bone resorption (postmenopausal state).
  \item Many patients were premenopausal at enrolment; benefit accumulates
        only after their transition to postmenopause (months 20--60).
  \item Sparse departure: hazard difference distributed across a
        mid-to-late window, null outside it.
\end{itemize}

\textbf{Test results}
\begin{itemize}
  \item Log-rank: $p = 0.305$ (NS).
  \item All four weighted alternatives: $p \geq 0.28$ (NS) — no fixed
        weight captures the mid-to-late window unaffected by the null
        early period.
  \item HCHG: $p = 0.005$; Fisher: $p = 0.004$
        (Figure~\ref{fig:main}C,G).
  \item Validated by published subgroup analyses: significant DFS benefit
        in postmenopausal patients \citep{coleman2011,coleman2014}.
\end{itemize}


\subsection*{Domain 4 — Corticosteroid in cirrhosis: time-varying effect
  (public IPD)}

\textbf{Trial and data}
\begin{itemize}
  \item CSL trial \citep{csl_data}: prednisone ($n=226$) vs placebo
        ($n=220$), histologically verified liver cirrhosis
        (Copenhagen, 1962--69). 270 deaths.
  \item Endpoint: overall survival. Follow-up up to $\sim$13 years.
  \item \textbf{Raw IPD genuinely public} --- distributed in the CRAN
        \texttt{timereg} package and the open SurvSet repository
        \citep{survset2022}; \emph{not} reconstructed from a KM curve.
  \item Note: source CSV is counting-process (long) format with a
        time-varying prothrombin covariate; collapsed to one row per
        patient (final follow-up time and status) for the two-sample test.
\end{itemize}

\textbf{Biological mechanism and survival pattern}
\begin{itemize}
  \item Corticosteroids have a time-varying effect in cirrhosis: early
        harm (immunosuppression, infection, fluid retention) followed by
        later anti-inflammatory benefit in actively inflamed subgroups.
  \item Survival curves cross more than once; treatment effect concentrated
        in a few non-contiguous windows (sign-changing).
  \item Canonical non-PH dataset, widely used to illustrate
        time-varying-coefficient models.
\end{itemize}

\textbf{Test results}
\begin{itemize}
  \item Log-rank: $p = 0.384$ (NS).
  \item All four weighted alternatives: $p \geq 0.13$ (NS) —
        Gehan-Wilcoxon $0.51$, Tarone-Ware $0.38$, Peto-Prentice $0.45$,
        Fleming-Harrington\,(1,1) $0.14$.
  \item HCHG: $p = 0.002$ (significant); Fisher: $p = 0.034$
        (Figure~\ref{fig:main}D,H).
  \item Demonstrates HCHG on genuinely public raw data, with a multi-window
        sign-changing departure no single weight can capture.
\end{itemize}


\textbf{Cross-domain summary}
\begin{itemize}
  \item HCHG achieves $p \leq 0.014$ in all four domains; log-rank is
        non-significant throughout ($p \geq 0.26$).
  \item GW ($p=0.041$) and PP ($p=0.049$) reach $p<0.05$ in CheckMate~057
        only — the sole weighted results below $0.05$; elsewhere all
        weighted tests $p \geq 0.13$. Their early emphasis coincidentally
        aligns with that trial's crossing signal and does not generalise.
  \item Fisher combination also significant in all four, consistent
        with sparse localised signals. MinP significant in CheckMate~057
        and AZURE.
  \item The four patterns differ (early+late crossing, transient-early,
        mid-to-late accumulated, multi-window sign-changing), yet HCHG
        detects all without any per-dataset tuning.
\end{itemize}


% =======================================================================
\section{Conclusion}
% =======================================================================

\begin{itemize}
  \item \texttt{lifelines-hc} provides a Python implementation of HCHG,
        the test specifically designed for sparse hazard departures at
        unknown temporal locations.
  \item Unlike weighted log-rank tests, HCHG requires no a priori
        specification of where the hazard departure occurs.
  \item Diagnostic output (\texttt{suspected\_deviations()}) identifies
        which time windows drive the signal.
  \item Fisher combination and MinP provided as complementary omnibus
        alternatives sharing the same hypergeometric infrastructure.
  \item Practical complement to standard survival testing whenever the
        proportional-hazards assumption is uncertain.
\end{itemize}


% =======================================================================
%  Back matter
% =======================================================================

\section*{Acknowledgements}
TODO.

\section*{Conflict of Interest}
None declared.

\section*{Funding}
TODO.

\section*{Data availability}
\begin{itemize}
  \item IPD for CheckMate~057, COMET-1, and AZURE reconstructed from
        published KM curves via \citet{guyot2012} as implemented in the
        \texttt{kmdata} R package (\url{https://github.com/raredd/kmdata}).
  \item CSL liver-cirrhosis IPD is fully public: CRAN \texttt{timereg}
        package (dataset \texttt{csl}) and the SurvSet repository
        \citep{survset2022} (\url{https://github.com/ErikinBC/SurvSet}).
  \item All analysis code:
        \url{https://github.com/alonkipnis/lifelines-hc/tree/main/AppNote}.
\end{itemize}


% =======================================================================
%  Figure  (4-domain version; adds CSL relative to main.tex)
% =======================================================================

\begin{figure*}[t!]
  \centering
  \includegraphics[width=\textwidth]{figure1_4dom}
  \caption{%
    \textbf{HCHG detects non-proportional hazard departures
    across four clinical datasets with distinct temporal patterns.}
    Top row (A--D): Kaplan-Meier curves with HCHG-flagged intervals
    (shaded). Bottom row (E--H): per-interval \emph{signed}
    $-\!\log_{10}$(hypergeometric p-value) --- upward bars mark excess
    events in the treatment arm, downward bars excess in the control arm.
    Dashed lines are the per-direction HCHG thresholds; bars reaching
    beyond either line are highlighted in red and coincide exactly with
    the intervals shaded above.
    %
    \textbf{(A,E)}~CheckMate~057 PFS (nivolumab vs.\ docetaxel, $n=582$);
    log-rank $p=0.351$, HCHG $p=0.001$. Curves cross at $\sim\!3$ months
    before nivolumab gains a durable late advantage.
    %
    \textbf{(B,F)}~COMET-1 OS (cabozantinib vs.\ prednisone, $n=1028$);
    log-rank $p=0.262$, HCHG $p=0.014$. Early OS advantage (months~3--7)
    during bone-lesion response, fading as resistance develops.
    %
    \textbf{(C,G)}~AZURE DFS (zoledronic acid vs.\ control, $n=3359$);
    log-rank $p=0.305$, HCHG $p=0.005$. Menopause-dependent benefit
    accumulates in months~20--60.
    %
    \textbf{(D,H)}~CSL OS (prednisone vs.\ placebo, liver cirrhosis,
    $n=446$); log-rank $p=0.384$, HCHG $p=0.002$. Time-varying
    corticosteroid effect with multiple curve crossings; \emph{fully
    public} raw IPD (CRAN \texttt{timereg} / SurvSet).
    In every dataset the log-rank test is non-significant while HCHG is
    significant. For the first three, IPD reconstructed
    via \citet{guyot2012}.%
  }
  \label{fig:main}
\end{figure*}


% =======================================================================
%  Bibliography  (unchanged from main.tex)
% =======================================================================

\bibliographystyle{plainnat}
\bibliography{refs}

\end{document}
